\documentclass[11pt]{article}
\usepackage[margin=1in]{geometry}
\usepackage{amsmath,amssymb,amsthm}
\usepackage[hidelinks]{hyperref}
\newcommand{\Tr}{\operatorname{Tr}}
\newcommand{\rank}{\operatorname{rank}}
\newcommand{\im}{\operatorname{Im}}
\newcommand{\norm}[1]{\left\|#1\right\|_{\mathrm F}}
\newtheorem{lemma}{Lemma}
\newtheorem{theorem}{Theorem}
\title{A trace inequality for four orthogonal projections\\
\large Complete solution of MathDB 332493}
\author{Principia Math in Progress}
\date{12 September 2026}
\begin{document}
\maketitle

\paragraph{Task 1. Solve the four-projection trace conjecture.}
Let $P,Q,R,S$ be orthogonal projections on $\mathbb C^K$ such that
\[
 PQ=0,\qquad RS=0,\qquad
 \im(P)\cap\im(R)=\{0\},\qquad
 \im(Q)\cap\im(S)=\{0\}.
\]
The conjecture additionally assumes
\[
 \rank P+\rank Q=\rank R+\rank S,\qquad
 \rank P+\rank R=\rank Q+\rank S.
\]
Define, using Moore--Penrose inverse square roots,
\[
\begin{aligned}
 P'&=(P+R)^{-1/2}P(P+R)^{-1/2},&
 R'&=(P+R)^{-1/2}R(P+R)^{-1/2},\\
 Q'&=(Q+S)^{-1/2}Q(Q+S)^{-1/2},&
 S'&=(Q+S)^{-1/2}S(Q+S)^{-1/2}.
\end{aligned}
\]
Banica and Nechita's Conjecture 6.6 asks whether
\begin{equation}\label{target}
 \Tr(P'Q')+\Tr(R'S')\leq \Tr(PR)+\Tr(QS).
\end{equation}

\begin{theorem}
Inequality \eqref{target} holds. In fact, the two rank identities are
unnecessary.
\end{theorem}

We first prove a scalar estimate. Every matrix norm below with subscript
$\mathrm F$ is the Frobenius norm; a contraction means a matrix of
operator norm at most one.

\begin{lemma}\label{scalar}
For $\lambda,\mu\in(-1,1)$, put
\[
 a=\frac1{\sqrt{(1+\lambda)(1+\mu)}},\qquad
 b=\frac1{\sqrt{(1-\lambda)(1-\mu)}}.
\]
Then $(a-b)^2\leq(\lambda^2+\mu^2)(a^2+b^2)$.
\end{lemma}
\begin{proof}
Set $t=\lambda^2+\mu^2$ and $w=\lambda\mu$. Since $1+w>0$,
\[
 \frac{(a-b)^2}{a^2+b^2}
 =1-\frac{\sqrt{1-t+w^2}}{1+w}.
\]
Thus it suffices to prove
\begin{equation}\label{scalar-root}
 \sqrt{1-t+w^2}\geq(1+w)(1-t).
\end{equation}
If $t\geq1$, the right side is nonpositive.
If $t<1$ and $w\leq0$, the left side is at least
$\sqrt{1-t}\geq1-t\geq(1+w)(1-t)$.
Finally, if $t<1$ and $w\geq0$, both sides of
\eqref{scalar-root} are nonnegative and the difference of their squares is
\[
 1-t+w^2-(1+w)^2(1-t)^2
 =(t-2w)(1-t)+wt\bigl(2(1-t)+w(2-t)\bigr)\geq0.
\]
Here $t-2w=(\lambda-\mu)^2\geq0$. This proves the lemma.
\end{proof}

\begin{lemma}\label{block}
Let $X=X^*$ and $Y=Y^*$ have spectra in $(-1,1)$, and suppose
\[
 \begin{pmatrix}I+X&C\\ C^*&I+Y\end{pmatrix}\geq0,\qquad
 \begin{pmatrix}I-X&-C\\ -C^*&I-Y\end{pmatrix}\geq0.
\]
Define
\[
 D_\pm=(I\pm X)^{-1/2}C(I\pm Y)^{-1/2}.
\]
Then
\begin{equation}\label{block-bound}
 \norm{\frac{D_+-D_-}{2}}^2
 \leq\frac12\bigl(\Tr X^2+\Tr Y^2\bigr).
\end{equation}
\end{lemma}
\begin{proof}
Congruence by the inverse square roots of the diagonal blocks gives
\[
 \begin{pmatrix}I&D_+\\D_+^*&I\end{pmatrix}\geq0,\qquad
 \begin{pmatrix}I&-D_-\\-D_-^*&I\end{pmatrix}\geq0.
\]
The Schur complement shows that both $D_+$ and $D_-$ are contractions.
Consequently, every row and every column of either matrix has squared
Euclidean norm at most one.

Choose separate orthonormal eigenbases for $X$ and $Y$, with eigenvalues
$\lambda_i$ and $\mu_j$. Write $c_{ij}$ for the entries of $C$ in these
bases. Entrywise, Lemma \ref{scalar} gives
\[
 |(D_+-D_-)_{ij}|^2
 \leq(\lambda_i^2+\mu_j^2)
       \bigl(|(D_+)_{ij}|^2+|(D_-)_{ij}|^2\bigr).
\]
Summing and using the row and column bounds yields
\[
\begin{aligned}
 \norm{\frac{D_+-D_-}{2}}^2
 &\leq\frac14\sum_{i,j}(\lambda_i^2+\mu_j^2)
       \bigl(|(D_+)_{ij}|^2+|(D_-)_{ij}|^2\bigr)\\
 &\leq\frac12\left(\sum_i\lambda_i^2+\sum_j\mu_j^2\right).
\end{aligned}
\]
This is \eqref{block-bound}.
\end{proof}

\begin{proof}[Proof of the theorem]
Choose matrices $p,q,r,s$ whose columns are orthonormal bases of
the ranges of $P,Q,R,S$, respectively. Thus $P=pp^*$ and similarly
for the other projections. Put
\[
 F=[\,p\ \ r\,],\qquad H=[\,q\ \ s\,],\qquad
 F^*F=I+X,\quad H^*H=I+Y,\quad C=F^*H.
\]
The zero intersection assumptions say exactly that $F$ and $H$ have
linearly independent columns. Hence $I+X$ and $I+Y$ are positive definite.
The Gram matrix
\[
 G=\begin{pmatrix}I+X&C\\C^*&I+Y\end{pmatrix}
   =[\,F\ \ H\,]^*[\,F\ \ H\,]
\]
is positive semidefinite.

Let
\[
 J=\operatorname{diag}(I_{\rank P},-I_{\rank R}),\qquad
 L=\operatorname{diag}(I_{\rank Q},-I_{\rank S}).
\]
The matrices $X$ and $Y$ have only off-diagonal blocks:
\[
 X=\begin{pmatrix}0&p^*r\\r^*p&0\end{pmatrix},\qquad
 Y=\begin{pmatrix}0&q^*s\\s^*q&0\end{pmatrix}.
\]
Also $p^*q=0$ and $r^*s=0$, so
\[
 C=\begin{pmatrix}0&p^*s\\r^*q&0\end{pmatrix}.
\]
It follows that $JXJ=-X$, $LYL=-Y$, and $JCL=-C$. In particular,
$I-X=J(I+X)J$ and $I-Y=L(I+Y)L$ are positive definite. Thus $X,Y$
have spectra in $(-1,1)$. Conjugating $G$ by
$\operatorname{diag}(J,L)$ supplies the second positive matrix in
Lemma \ref{block}.

Set
\[
 U=F(I+X)^{-1/2}=[\,u_P\ \ u_R\,],\qquad
 V=H(I+Y)^{-1/2}=[\,v_Q\ \ v_S\,].
\]
These matrices have orthonormal columns. The singular value decomposition
of $F$ gives
\[
 (FF^*)^{-1/2}F=F(F^*F)^{-1/2};
\]
the inverse on the left is the Moore--Penrose inverse square root.
Since $FF^*=P+R$, the projections onto the respective column groups of
$U$ are exactly $P',R'$. Likewise, the groups of $V$ give $Q',S'$.

We have $U^*V=D_+$, and functional calculus together with $JCL=-C$
gives
\[
 JD_+L
 =(I-X)^{-1/2}(JCL)(I-Y)^{-1/2}
 =-D_-.
\]
The matrix
\[
 \frac{D_++JD_+L}{2}=\frac{D_+-D_-}{2}
\]
retains exactly the two diagonal rectangular blocks $u_P^*v_Q$ and
$u_R^*v_S$ of $D_+$. Therefore
\[
\begin{aligned}
 \Tr(P'Q')+\Tr(R'S')
 &=\norm{u_P^*v_Q}^2+\norm{u_R^*v_S}^2\\
 &=\norm{\frac{D_+-D_-}{2}}^2\\
 &\leq\frac12(\Tr X^2+\Tr Y^2)\\
 &=\norm{p^*r}^2+\norm{q^*s}^2\\
 &=\Tr(PR)+\Tr(QS).
\end{aligned}
\]
No rank identity was used. Zero-dimensional column groups are allowed;
if an entire frame has no columns, all the corresponding sums vanish.
\end{proof}

\paragraph{Verification and scope.}
The argument proves the conjecture for all finite dimensions over
$\mathbb C$, including zero and unequal ranks. The accompanying script
checks the scalar polynomial certificate exactly and checks the matrix
identities on 724 reproducible real and complex examples, including
unequal ranks and nearly intersecting ranges. These numerical checks
are supplementary; the proof above is exact. The higher-$N$
orthogonalization and convergence conjectures are not asserted here.

\paragraph{Task 2. Determine the equality cases.}
Use the equality conditions in the scalar and contraction estimates
to characterize all quadruples for which \eqref{target} is an equality.

\paragraph{Task 3. Quantify the decrease.}
Derive explicit improvements to \eqref{target} when the nonzero principal
angles between the relevant projection ranges obey prescribed bounds.

\paragraph{Task 4. Extend the block estimate.}
Determine which Schatten-norm or trace-class versions of Lemma
\ref{block} hold, and state the hypotheses needed in infinite dimension.

\paragraph{Task 5. Investigate larger orthogonalization arrays.}
Identify an analogue of the two-block grading argument for three or
more projection groups and test its consequences for the higher-$N$
inequalities in the source.

\begin{thebibliography}{9}
\bibitem{BN}
T. Banica and I. Nechita,
\emph{Flat matrix models for quantum permutation groups},
Advances in Applied Mathematics \textbf{83} (2017), 24--46.
\href{https://doi.org/10.1016/j.aam.2016.09.001}{doi:10.1016/j.aam.2016.09.001}.
Conjecture 6.6, p.~20 of the
\href{https://arxiv.org/abs/1602.04456}{arXiv version} and the author's
currently linked manuscript.
\end{thebibliography}
\end{document}
